% ---------------------------------------------------------------------
\subsection{Partitioned convolution (\fn{PartConv.h}, \fn{FastFFT.h})}
\label{sec:math-partconv}
Every long filter in Rysum (the impulse-response room, the room tail of the \ui{Saal}, the binaural and transaural
filters) is a convolution $y = x * h$ with a response of up to ten seconds, i.e.\ up to $4.8\cdot10^5$ taps at
\SI{48}{kHz}. Done directly that is about $10^{11}$ multiplications per second and channel; done with one FFT of the
whole length it has a latency of ten seconds. \fn{PartitionedConvolution} is the usual way between the two
\citep{gardner1995}: uniformly partitioned overlap-save, with the work of each block spread over time.

\paragraph{Uniformly partitioned overlap-save.} The response is cut into $P$ partitions of $B$ samples,
$h_q[n] = h[qB+n]$, $n = 0\dots B-1$, each zero-padded to $N = 2B$ and transformed once when the room is loaded:
$H_q = \mathcal F_N\{h_q, 0_B\}$ (\fn{RealFFT::forwardScaled}, which includes the factor $1/N$ of the inverse
transform). The input is collected in blocks $x_j$ of $B$ samples; at the end of each block the window of the last two
blocks is transformed, $X_j = \mathcal F_N\{x_{j-1}, x_j\}$, and stored in a \emph{frequency-domain delay line} (FDL)
that keeps the last $P$ spectra. The output block is
\begin{equation}
  y_j = \text{last $B$ samples of}\ \mathcal F_N^{-1}\Bigl\{\sum_{q=0}^{P-1} X_{j-q}\cdot H_q\Bigr\}.
\end{equation}
The product of two spectra of length $2B$ is a \emph{circular} convolution of length $2B$; since $h_q$ has only $B$
non-zero samples, the circular wrap-around falls entirely into the first $B$ output samples, and the last $B$ are the
exact linear convolution of the window with $h_q$ (this is why it is called overlap-\emph{save}: the first half is
thrown away). Summing over $q$ with the delayed spectra $X_{j-q}$ adds the contribution of each partition $qB$
samples later. The cost per block is one forward FFT per input, one inverse FFT per output and $P$ complex
multiply-adds of $B+1$ bins per input--output pair (\fn{multiplyAdd}, PFFFT's \code{zconvolve\_accumulate}
\citep{pffft}, in PFFFT's internal bin order, which is never converted because nothing but products and sums is done
in the frequency domain).

\paragraph{Spreading the work.} If all $P$ products were done when a block is complete, one audio callback in every
$B/b$ (for a host block $b$) would carry the whole load: a CPU spike that can cause a dropout even though the
average load is small. The class therefore counts the work of the coming block in \emph{units} (one partition
product; an inverse FFT counts as 4) and, after every \fn{push} of $n$ samples, does as many as are due,
\begin{equation}
  u_\mathrm{done} \ge \Bigl\lceil U\cdot\frac{p}{B}\Bigr\rceil ,
\end{equation}
where $p$ is the position in the block and $U$ the total. Only products whose input spectra already exist can be
done in advance: $X_{j-q}$ for $q\ge1$ is known during block $j$, $X_j$ only at its end.
\begin{itemize}[nosep]
  \item \emph{Lag 1}: all products with $q\ge 1$ are spread over block $j$; at the boundary the product with the
    newest spectrum ($q=0$) is added and the inverse FFT done, and the result is played during block $j+1$. The output
    is the response delayed by one block.
  \item \emph{Lag 2}: the stage is used for a part of the response that is silent for its first $2B$ samples, so even
    the newest partition it needs is a block old: everything, including the inverse FFT, is spread over the block
    before, and the boundary only swaps buffers.
\end{itemize}
Changes of the response (\fn{setTerms}) take effect at the next block boundary, never in the middle of a sum, so a
room can be swapped without a half-and-half block. Each output is a sum of \emph{terms} (input $k$ convolved with
partitions $q_0\dots q_0+Q-1$ of a response, with an optional sign), which is what lets one stage compute the
rotation sums of the next subsection.

\paragraph{Non-uniform partitions.} Small partitions keep the latency low but make many products; large partitions
make few products but need a long wait. The impulse-response room uses three stages in parallel, each a
\fn{PartitionedConvolution} fed with the same inputs:
\begin{center}
\begin{tabular}{@{}lllll@{}}
\toprule
stage & $B$ & lag & covers (samples of the response, after the delay of one block) & partitions \\
\midrule
A & 512 & 1 & $[0,\ 3584)$ & 7 \\
B & 2048 & 2 & $[3584,\ 15872)$ & 6 \\
C & 8192 & 2 & $[15872,\ \infty)$ & as many as needed \\
\bottomrule
\end{tabular}
\end{center}
The boundaries follow from the lags: a lag-2 stage with block $b$ can start where the response delayed by one
stage-A block has been silent for $2b$ samples, $2b - 512$, i.e.\ $3584 = 2\cdot2048-512$ and $15872 = 2\cdot 8192 - 512$.
For a response of \SI{10}{s} ($4.8\cdot10^5$ samples) uniform partitions of 512 need 938 products per block of 512
samples; the three stages need $7+6+57 = 70$ partitions. A product in stage B or C is four or sixteen times larger
than one in stage A but is also done only once every four or sixteen blocks, so per sample each partition costs about
the same, and the work drops from 938 to about 70 units (\S\ref{sec:perf} gives the counts for a \SI{3}{s} response
and the measurements). What is paid for it is the larger FFTs of the later stages, which cost
$O(\log B)$ more per sample, and the memory for their spectra. The three stages run in lock-step: \fn{Convolver::process} always advances by the
largest number of samples that brings no stage past its own block boundary.

% ---------------------------------------------------------------------
\subsection{The impulse-response room (\fn{Convolver.cpp})}
\label{sec:math-ir}
This is the function behind \S\ref{sec:ir}. It turns a recorded response, mono, stereo or Ambisonic, into a room in
which every pipe's early reflections come from that pipe's own direction, without a convolution per pipe.

\paragraph{Reading the file (\fn{Convolver::create}).} A file with $(L+1)^2$ channels is read as AmbiX of order $L$
(at most 5); any other number of channels as mono or stereo (more than two: the first two). First-order FuMa
($W,X,Y,Z$) becomes AmbiX by reordering to $W,Y,Z,X$ and multiplying $W$ by $\sqrt2$ (FuMa stores $W$ at
$-\SI{3}{dB}$); N3D becomes SN3D by dividing degree $l$ by $\sqrt{2l+1}$. The response is resampled to the host rate
with a four-point Lagrange interpolator, trailing chunks of 256 samples below $10^{-4}$ of the peak
($\SI{-80}{dB}$) are cut, and the response is scaled so that the energy of its omnidirectional channel is one,
$\sum_n w[n]^2 = 1$: rooms of different loudness then have the same reverb level for the same send. A response that
has died away ends with a linear fade of at most \SI{50}{ms}; one that still sounds where the ten-second limit cut it
ends with a fade $\bigl(\tfrac12+\tfrac12\cos\pi t\bigr)^2$ of up to \SI{2}{s}, which has no corner at its start and
falls ever faster, so the cut is not heard as one.

\paragraph{The direct sound and its direction.} The direct sound is taken to be the largest $|w[n]|$, at $t_d$. For
an Ambisonic response its direction is read from the first-order intensity over \SI{-0.5}{}\dots\SI{+2}{ms} around it:
\begin{equation}
  \varphi_0 = \operatorname{atan2}\Bigl(\sum_n w[n]\,Y[n],\ \sum_n w[n]\,X[n]\Bigr)
\end{equation}
(ACN channel 1 is $Y$, channel 3 is $X$). A plane wave from $\varphi_s$ is encoded as $W = s$, $X = s\cos\varphi_s$,
$Y = s\sin\varphi_s$, so the two sums are $E\cos\varphi_s$ and $E\sin\varphi_s$ with the energy $E$ of the window, and
their angle is the direction the direct sound came from (the two sums are the first-order intensity vector, with the
sign of the encoding convention). Reflections inside the short window bias it only slightly, as they are weaker. The whole response is then turned by
$-\varphi_0$ (\fn{rotateYaw}, below), so its direct sound points to the front. A mono or stereo response (the omni
signal $w$ being the file or the mid of the stereo pair) has no direction; its early part is encoded as a point source
in front, $A_c[n] = w[n]\,Y_c(0,0)$, up to the output order. Around $t_d$ (from \SI{-1}{} to \SI{+3}{ms}, with
raised-cosine edges of \SI{0.5}{ms}) the response is scaled by \ui{Direktschall der IR} (default $0.2$): the dry pipes
are already the direct sound, and a second copy of it from the loudspeaker's direction would pull every pipe there.

\paragraph{Early and late.} The response is split \SI{80}{ms} after the direct sound, rounded to a partition of 512,
with a quarter-cosine crossfade over one partition: the early part is multiplied by $c(n) = \cos(\tfrac\pi2 t)$ and
the late part by $1-c(n)$. For the omni channel both parts contain the same samples, so the two weights sum to one
and $W$ is reproduced exactly; the split only decides which samples are rotated per pipe and which are shared.

\paragraph{Rotation about the vertical axis (\fn{rotateYaw}).} Turning a sound field by an angle $\alpha$ about the
vertical axis moves every direction from $\varphi$ to $\varphi+\alpha$. The elevation part $N_l^mP_l^m(\sin\vartheta)$
of each harmonic does not change, and the azimuth part only mixes the pair $(m,-m)$:
\begin{equation}
  \begin{pmatrix} a'_{l,m} \\ a'_{l,-m}\end{pmatrix} =
  \begin{pmatrix} \cos m\alpha & -\sin m\alpha \\ \sin m\alpha & \cos m\alpha \end{pmatrix}
  \begin{pmatrix} a_{l,m} \\ a_{l,-m}\end{pmatrix},\qquad m = 1\dots l,
\end{equation}
which follows from $\cos m(\varphi-\alpha) = \cos m\varphi\cos m\alpha + \sin m\varphi\sin m\alpha$ and the same for
the sine: a field that \emph{contained} a source at $\varphi_s$ now contains it at $\varphi_s+\alpha$. Channels with
$m=0$ are unchanged. The cost is a $2\times2$ rotation per channel pair and sample, with no interpolation and no loss.

\paragraph{The rotation per pipe, moved before the convolution.} What should come out of the room for pipe $p$ at
azimuth $\varphi_p$ with signal $s_p$ is its signal convolved with the early response turned by $\varphi_p$. Summed over
all pipes, channel $(l,m)$ with $m>0$ of the early part is
\begin{align}
  e_{l,m} &= \sum_p \bigl(\cos(m\varphi_p)\,h_{l,m} - \sin(m\varphi_p)\,h_{l,-m}\bigr) * s_p
           = h_{l,m} * C_m - h_{l,-m} * S_m, \\
  e_{l,-m} &= h_{l,m} * S_m + h_{l,-m} * C_m,\qquad e_{l,0} = h_{l,0} * C_0,
\end{align}
with the \emph{circular-harmonic sends}
\begin{equation}
  C_0 = \sum_p s_p,\qquad C_m = \sum_p \cos(m\varphi_p)\,s_p,\qquad S_m = \sum_p \sin(m\varphi_p)\,s_p,\qquad m = 1\dots L.
\end{equation}
Because convolution is linear, the cosines and sines, which depend on the pipe, can be applied to the pipe's
signal before the convolution, and the responses, which do not, after. Instead of one convolution per pipe there are
$2L+1$ sends; each output channel with $m\neq0$ is the sum of two terms and each with $m=0$ of one, which is exactly
the term structure of \fn{PartitionedConvolution} (the minus sign is a term with \code{negate}). The engine builds
the sends in \fn{renderRange} by the angle-addition recurrence for $\cos m\varphi_p$, $\sin m\varphi_p$, with the gains
ramped across the block like every other gain; the input order is $C_0, C_1, S_1, C_2, S_2,\dots$. With \ui{folgt
Pfeifen} off, every $\varphi_p$ is the same fixed angle and the room is simply the measured one, turned.

Two properties follow from the construction. The rotation is about the vertical axis only: a pipe above the listener
gets the early reflections of a pipe on the horizon at its azimuth (a rotation that also tilts would mix all
$2l+1$ channels of a degree and need $(L+1)^2$ sends). And for a mono response the early part is the omni response
placed at each pipe's direction: each pipe is heard with its own coloured, delayed copy coming from where it stands,
but the reflections are not spread around it, because a mono file contains no directions to spread.

\paragraph{The late part (\fn{diffuseCopy}).} The late part is fed by $C_0$ alone: after \SI{80}{ms} a real room's
field is nearly diffuse and no longer depends much on where the source stood. Its omni channel is the measured
tail $w_\mathrm{late}$. Every other channel $c$ of degree $l$ is
\begin{equation}
  h_c^\mathrm{late} = \sqrt{1-\varepsilon}\;h_c^\mathrm{measured} \;+\; \sqrt{\varepsilon}\;
  \sqrt{\frac{E_W}{(2l+1)\,E_{d_c}}}\;d_c ,
\end{equation}
where $\varepsilon$ is \ui{Umhüllung}, $E_W$ the energy of the omni tail, the first term exists only for channels the
file has, and $d_c$ is a \emph{random-phase copy} of the omni tail made with its own seed per channel:
\begin{enumerate}[nosep]
  \item cut $w_\mathrm{late}$ into frames of 1024 samples with a hop of 256, each multiplied by a Hann window;
  \item transform each frame, keep every bin's magnitude $|W_k|$ and replace its phase by a uniformly random one;
  \item transform back, multiply by the Hann window again and overlap-add; a squared Hann window at 75\,\% overlap
    sums to $1.5$, hence the division by $1.5$.
\end{enumerate}
The copy has the same spectrogram as the measured tail (the same decay at every frequency, the same colour) but
its fine structure is new, so copies with different seeds are mutually uncorrelated. Uncorrelated channels with
equal statistics are what a diffuse field looks like in the spherical-harmonic domain, and in SN3D the expected
energy of a channel of degree $l$ in an isotropic field is $1/(2l+1)$ of the omni channel's, which is the factor in
the formula. The blend uses square roots so that it preserves energy ($(\sqrt{1-\varepsilon})^2+(\sqrt\varepsilon)^2 = 1$
for uncorrelated parts). This is also what lifts a mono, stereo or first-order response to the full fifth order: the
missing channels are synthesised, the measured ones kept as far as \ui{Umhüllung} allows.

\paragraph{Stereo output.} For a stereo output the room is computed at first order and the field is turned into two
virtual cardioid microphones pointing left and right, $L = \tfrac12(W+Y)$, $R = \tfrac12(W-Y)$; a stereo response on
a stereo output is played as recorded instead, left and right as they are. The whole wet signal is delayed by one
block of 512 samples (\SI{10.7}{ms} at \SI{48}{kHz}) by the lag-1 stage; since the response's own direct sound is
turned down and the first reflections of a church arrive later than that, the delay is not audible as such.
